\paragraph{\trEquivariant\trTwistedComplex。}\label{sub:equivariant-twisted-complexes}
作为预备的一步，我们需要讨论\trAdditive{}扩大的\trEquivariant{}版本
$\Sigma^{\mathrm{eq}}\scr{A}$。$\Sigma^{\mathrm{eq}}\scr{A}$ 的\trObject{}是
$\Sigma\scr{A}$ 的一个\trObject，亦即一个如\eqref{eq:sigma-formal-direct-sum}
中那样写成形式\trDirectSum{}的三元组 $(I,\{X^i\},\{V^i\})$，连同以下附加数据：
标号集 $I$ 上的一个\trInvolution{} $i\mapsto\underline{i}$，使得
$X^{\underline{i}}=\epsilon_{\scr{A}}(X^i)$；以及诸\trVectorSpace{}之间的%
\trGraded{}映射
$\epsilon_{V^i}\colon V^i\to V^{\underline{i}}$，使得对一切 $i$ 都有
$\epsilon_{V^{\underline{i}}}\circ\epsilon_{V^i}=\mathrm{Id}$。\trMorphism{}空间与%
\trComposition{}映射的定义与非\trEquivariant{}版本相同。每个\trMorphism{}空间都
带有一个诱导的\trInvolution；按\eqref{eq:sigma-morphism-spaces} 的记号，它是取
\begin{equation*}
\epsilon_{V_1^j}\circ-\circ\epsilon_{V_0^{\underline{i}}}\colon
\hom_{\KK}(V_0^i,V_1^j)\longrightarrow
\hom_{\KK}\bigl(V_0^{\underline{i}},V_1^{\underline{j}}\bigr)
\end{equation*}
与
\begin{equation*}
\epsilon_{\scr{A}}\colon
\hom_{\scr{A}}\bigl(X_0^i,X_1^j\bigr)\longrightarrow
\hom_{\scr{A}}\bigl(X_0^{\underline{i}},X_1^{\underline{j}}\bigr)
\end{equation*}
的\trTensorProduct{}得到的。\textbf{\trEquivariant\trTwistedComplex}（equivariant
twisted complex）是一个对 $(X,\delta_X)$，由
$X\in\Ob\Sigma^{\mathrm{eq}}\scr{A}$ 与一个位于不变部分
$\hom^1_{\Sigma^{\mathrm{eq}}\scr{A}}(X,X)^{\ZZ/2}$ 中的\trDifferential{}%
$\delta_X$ 组成，并带有通常的条件（存在一个\trFiltration，现在应当是按%
\trEquivariant{}子\trComplex{} $F^iX$ 的\trFiltration；以及%
\trMaurerCartan{}方程）。%
\trMorphism{}空间承袭自 $\Sigma^{\mathrm{eq}}\scr{A}$，而\trComposition{}映射一如
\eqref{eq:twisted-composition}，其结果是一个\trACategory{}%
$Tw^{\mathrm{eq}}\scr{A}$，其上带有一个在\trObject{}上是平凡的朴素 $\ZZ/2$-作用。

最简单的\trEquivariant\trTwistedComplex{}莫过于 $\scr{A}$ 的一个单独的
$\epsilon_{\scr{A}}$-不动\trObject{} $X$（更精确地说，取一个不动的一元素集
$I=\{*\}$、一个一维\trVectorSpace{} $V^*=\KK$ 连带平凡映射
$\epsilon_{V^*}=\mathrm{Id}$，以及 $X^*=X$）。一个不那么平凡的例子，是 $\scr{A}$
的任意\trObject{}与其像的\trDirectSum{} $X\oplus\epsilon_{\scr{A}}X$；此时%
\trDifferential{}依然是平凡的，但\trInvolution{} $i\mapsto\underline{i}$ 交换两个%
\trDirectSummand。注意，第一类\trObject{}给出 $\scr{A}^{\mathrm{fix}}$ 到
$Tw^{\mathrm{eq}}\scr{A}$ 之中的一个\trFull{}\trEmbedding，且与诸 $\ZZ/2$-作用%
\trCompatible。推而广之，$\scr{A}^{\mathrm{inv}}$ 中的任何\trTwistedComplex{}都可以看作一个%
\trEquivariant\trTwistedComplex（集合 $I$ 与诸\trVectorSpace{} $V^i$ 上带平凡的%
\trInvolution），这样一来，$Tw(\scr{A}^{\mathrm{inv}})$ 就成为
$(Tw^{\mathrm{eq}}\scr{A})^{\mathrm{inv}}$ 的一个\trFull{} A∞-子\trCategory。
最后，我们还可以把这些\trInvolution{}统统忘掉，这就定义了一个\trFullyFaithful{}的%
\trAFunctor{} $Tw^{\mathrm{eq}}\scr{A}\to Tw\scr{A}$。

我们提几种常见运算的\trEquivariant{}版本。\trDirectSum{}可以按显然的方式定义。
其次，给定一个带 $\ZZ/2$-作用 $\epsilon_Z$ 的有限维\trGraded\trVectorSpace{} $Z$
与一个\trEquivariant\trTwistedComplex{} $X$，可以作出\trEquivariant\trTensorProduct{}%
$Y=Z\otimes X$：取诸\trTensorProduct{} $Z\otimes V^i$ 并配上联合的 $\ZZ/2$-作用
$\epsilon_Z\otimes\epsilon_{V^i}$，以及\trDifferential{}%
$\delta_Y=\id_Z\otimes\delta_X$。这包括了\trSuspension{}运算 $S$，它按通常的方式
定义为取 $Z=\KK[1]$ 并配平凡 $\ZZ/2$-作用。另一个有趣的情形是 $\ZZ/2$ 的非平凡%
\trCharacter{}的情形，即 $Z=\KK$ 但 $\epsilon_Z=-\id_Z$。此时，\trTensorProduct{}只不过把
诸 $\epsilon_{V^i}$ 的正负号全部反转，我们把这一过程记作 $X\mapsto\bar{X}$。一般的%
\trTensorProduct{} $Z\otimes X$ 本质上是 $X$ 与 $\bar{X}$ 的若干\trSuspension{}副本
的有限和，$Z$ 的每个不变（反不变）\trBasis{}元对应一个副本。最后，若
$c\in\hom_{Tw\scr{A}}(Y_0,Y_1)$ 是一个不变的零次\trCocycle，那么经由给
$SY_0\oplus Y_1$ 配上\trDifferential{} \eqref{eq:mapping-cone-twisted-complex}
而定义的\trMappingCone{}仍是一个\trEquivariant\trTwistedComplex。这表明
$(Tw^{\mathrm{eq}}\scr{A})^{\mathrm{inv}}$ 是\trTriangulated{}的
（$Tw^{\mathrm{eq}}\scr{A}$ 本身则不是）。

\begin{lemma}\label{thm:equivariant-cone-isomorphism}
设 $Y$、$Y_0$、$Y_1\in\Ob\scr{A}$ 都是 $\epsilon_{\scr{A}}$-不动的，并附加假定：
$Y$ 是自由的，且 $\Hom_{H(\scr{A})}(Y_0,Y_1)$ 上的 $\ZZ/2$-作用是平凡的。设在
$D(\scr{A})=H^0(Tw\scr{A})$ 中，$Y$ \trIsomorphic{}于某个\trMorphism{} $Y_0\to Y_1$ 的%
\trCone。那么，同样的断言在 $D(\scr{A}^{\mathrm{inv}})$ 中也成立。
\end{lemma}

\begin{proof}
按假定，$\Hom_{H(\scr{A})}(Y_0,Y_1)$ 是 $\ZZ/2$-不变的，因此在不改变\trCone{}的%
\trIsomorphism{}类的前提下，可以假设相关的\trCocycle{}%
$c\in\hom^0_{\scr{A}}(Y_0,Y_1)$ 本身是 $\ZZ/2$-不变的。在这种情形下，我们可以把
$\mathrm{Cone}(c)$ 作成 $Tw^{\mathrm{eq}}\scr{A}$ 的一个\trObject。实际上，它位于%
\trEmbedding{} $Tw(\scr{A}^{\mathrm{inv}})\subset Tw^{\mathrm{eq}}\scr{A}$ 的像
之中。由于遗忘\trFunctor{} $H^0(Tw^{\mathrm{eq}}\scr{A})\to H^0(Tw\scr{A})$ 是%
\trFullyFaithful{}的，我们便得到 $H^0(Tw^{\mathrm{eq}}\scr{A})$ 中的一个%
\trIsomorphism{} $Y\cong\mathrm{Cone}(c)$。接下来，因为 $Y$ 是自由的，可以利用%
\cref{thm:free-fixed-invariant-isomorphic}把这个\trIsomorphism{}做成不变的，
也就是说，它来自
$H^0(Tw(\scr{A}^{\mathrm{inv}}))$ 中的一个\trIsomorphism。
\end{proof}

\begin{lemma}\label{thm:equivariant-split-generation}
设 $\scr{A}$ 是\trCfinite{}的。设 $X$ 是一个 $\epsilon_{\scr{A}}$-不动\trObject。
设 $X$ \trQuasiRepresent{}某个\trObject{} $Y$ 及其 $\epsilon_{\scr{A}}$-像的%
\trDirectSum{} $Y\oplus\epsilon_{\scr{A}}Y$，而且还满足
$\Hom_{H(\scr{A})}(Y,\epsilon_{\scr{A}}Y)=0$。再设在 $D(\scr{A})$ 中，我们有%
\trIsomorphism
\begin{equation}\label{eq:iterated-twist-split}
T_{Y_1}\ldots T_{Y_m}(Y)\cong S(\epsilon_{\scr{A}}Y).
\end{equation}
这里诸 $Y_1,\ldots,Y_m$ 是 $\scr{A}$ 的 $\epsilon_{\scr{A}}$-不动\trObject，其
性质是对一切 $j<k$，$\Hom_{H(\scr{A})}(Y_j,Y_k)$ 上的 $\ZZ/2$-作用都是平凡的；
而 $S$ 是\trSuspension。那么 $X$ 位于 $Tw(\scr{A}^{\mathrm{inv}})$ 中由诸 $Y_k$
生成的\trTriangulated{}子\trCategory{}之内。
\end{lemma}

\begin{proof}
最透明的情形，是 $\scr{A}$ \trStrictlyUnital{}且 $\mu^1_{\scr{A}}=0$ 的情形（这
特别蕴含 $\hom_{\scr{A}}$ 空间本身必须是有限维的）。通过迭代
\eqref{eq:twist-cone-representative}，可以发现\eqref{eq:iterated-twist-split}
左端的\trObject{} $X$ 具有形式
\begin{equation}\label{eq:iterated-twist-expansion}
\begin{aligned}
C=Y\oplus\Bigl(\bigoplus \hom_{\scr{A}}(Y_{j_r},Y)[1]\otimes
\hom_{\scr{A}}(Y_{j_{r-1}},Y_{j_r})[1]\otimes\cdots\\
\cdots\otimes\hom_{\scr{A}}(Y_{j_1},Y_{j_2})[1]\otimes Y_{j_1}\Bigr),
\end{aligned}
\end{equation}
其中 $\bigoplus$ 取遍一切 $r\ge 1$ 与 $j_1<j_2\cdots<j_r$。与其把整个%
\trDifferential{} $\delta_C$ 显式写出，我们其实只需知道它的一般形状，而这容易从
构造中推出。其一，
\begin{itemize}
\item $\delta_C$ 严格缩减诸\trTensorProduct{}表达式的长度。特别地，它在第一个%
\trDirectSummand{} $Y$ 上为零。
\end{itemize}
这意味着 $C$ 可以写成某个\trMorphism{} $U\to Y$ 的\trCone，其中
$(U,\delta_U)\in\Ob Tw\scr{A}$ 包含\eqref{eq:iterated-twist-expansion} 中除第一个
之外的全部项，连同\trDifferential{}的相应部分。其二，
\begin{itemize}
\item $\delta_U$ 与 $C$ 的如下\trFiltration{}\trCompatible：只考虑使 $j_1\ge j$
（对某个 $j$）成立的那些\trDirectSummand。此外，相对于这一\trFiltration{}的对角项
（即相伴\trGraded{}空间上的诱导\trDifferential）是 $\partial_{Z_j}\otimes e_{Y_j}$，
其中 $\partial_{Z_j}$
是 $Z_j=\bigoplus\hom_{\scr{A}}(Y_{j_{r-1}},Y_{j_r})\otimes\cdots\otimes
\hom_{\scr{A}}(Y_j,Y_{j_2})$ 上的某个\trDifferential。这里，和取遍一切 $r$ 以及
满足 $j_2\ge j$ 的 $j_2<\cdots<j_r$（再细看便会发现 $\partial_{Z_j}$ 是 bar
\trDifferential{}的一种形式，但我们不会用到这一事实）。
\end{itemize}
这一点，连同 $\hom_{\scr{A}}(Y_j,Y_k)$ 在 $j<k$ 时整体 $\ZZ/2$-不变的假定，蕴含
$\delta_U$ 的每个\trCoefficients{}都是某个\trVectorSpace{}\trHomomorphism{}与
$\scr{A}$ 中一个 $\epsilon_{\scr{A}}$-不变\trMorphism{}的\trTensorProduct。因此，
$U$ 位于子\trCategory{} $Tw(\scr{A}^{\mathrm{inv}})\subset Tw\scr{A}$ 之中；而再
看一眼上面的陈述即可知道，它实际上位于 $Tw(\scr{A}^{\mathrm{inv}})$ 中由诸 $Y_k$
生成的\trTriangulated{}子\trCategory{}之内。鉴于\eqref{eq:iterated-twist-split}，
我们在 $H(Tw\scr{A})$ 中有如下形式的\trExactTriangle：
\begin{equation*}
\begin{tikzcd}
U \ar[rr] && Y \ar[dl] \\
& S(\epsilon_{\scr{A}}Y). \ar[ul,"{[1]}"] &
\end{tikzcd}
\end{equation*}
按假定，在 $H(\scr{A})$ 中从 $Y$ 到 $\epsilon_{\scr{A}}Y$ 没有任何\trDegree{}的
非平凡\trMorphism，因此该\trExactTriangle{}\trSplitting，也就是说
\begin{equation}\label{eq:splitting-free-middle}
X\cong Y\oplus\epsilon_{\scr{A}}Y\cong U
\end{equation}
在 $H^0(Tw\scr{A})$ 中成立。\eqref{eq:splitting-free-middle} 中的所有\trObject{}%
都是\trEquivariant\trTwistedComplex，而中间那个是自由的。于是，%
\cref{thm:equivariant-cone-isomorphism}之证明收尾处的论证便给出所要的
$H^0(Tw(\scr{A}^{\mathrm{inv}}))$ 中的\trIsomorphism。

要把结果推广到一般的 $\scr{A}$、去掉上面引入的技术性限制，做法如下。选取如%
\cref{thm:perturbation-typical-application}中的数据，但要注意使诸\trChainComplex{}%
$\hom_{\scr{A}}(X_0,X_1)$ 上的\trSplitting{}与\trHomotopy{}同诸 $\ZZ/2$-作用%
\trCompatible。唯一可能
引起潜在困难的场合，是 $X_0$ 与 $X_1$ 都是 $\epsilon_{\scr{A}}$-不变的情形；但在
那些场合，只需把这个\trComplex{}\trSplitting{}成不变与反不变的\trDirectSummand，
对每一部分
分别进行，最后再把结果拼在一起。\cref{thm:perturbation-lemma}之后的公式
\eqref{eq:perturbation-recursive-formulae} 保持既有的对称性，这意味着我们得到一个%
\trEquivariant\trQuasiisomorphism{} $\mathscr{F}\colon\scr{A}\to\scr{B}$，其中
$\scr{B}$ 是一个带朴素 $\ZZ/2$-作用、且 $\mu^1_{\scr{B}}=0$ 的 \trACategory。
接下来，我们要像\cref{thm:strictly-unital-modification}中那样使 $\scr{B}$ 成为%
\trStrictlyUnital{}的。那条引理之证明的第一步是平凡的：由于 $\mu^1_{\scr{B}}=0$，
\trCohomology{}幺元的代表元唯一（而且自动 $\ZZ/2$-不变），因此可以取 $\Phi^2=0$。
构造的其余部分不涉及任何选取，见\eqref{eq:formal-diffeomorphism-recursion}，因而
自然与我们的\trInvolution{}\trCompatible，这就给出一个\trEquivariant\trFormalDiffeomorphism{}%
$\Phi$，使得 $\widetilde{\scr{B}}=\Phi_{*}\scr{B}$ 是\trStrictlyUnital{}的。取
$\mathscr{G}=\Phi\circ\mathscr{F}$，并考察 $X$、$Y$ 与诸 $Y_k$ 在其下的像。显然
\begin{equation*}
\begin{aligned}
&\mathscr{G}X\cong\mathscr{G}Y\oplus\mathscr{G}(\epsilon_{\scr{A}}Y)
\cong\mathscr{G}Y\oplus\epsilon_{\widetilde{\scr{B}}}(\mathscr{G}Y)
\quad\text{在 }H^0(\widetilde{\scr{B}})\text{ 中},\\
&\Hom^{*}_{H(\widetilde{\scr{B}})}\bigl(\mathscr{G}Y,
\epsilon_{\widetilde{\scr{B}}}(\mathscr{G}Y)\bigr)=0,\\
&\Hom^{*}_{H(\widetilde{\scr{B}})}(\mathscr{G}Y_i,\mathscr{G}Y_j)
\text{ 是 }\ZZ/2\text{-不变的，对 }i<j.
\end{aligned}
\end{equation*}
据\cref{thm:twist-functorial-isomorphism}，\eqref{eq:iterated-twist-split} 的类似
命题也成立，于是我们正处于上面考虑过的特殊情形，也就是说，$\mathscr{G}X$ 位于
$Tw(\widetilde{\scr{B}}^{\mathrm{inv}})$ 中由诸 $\mathscr{G}Y_k$ 生成的%
\trTriangulated{}子\trCategory{}之内。但由于不变部分
$\mathscr{G}^{\mathrm{inv}}\colon\scr{A}^{\mathrm{inv}}\to
\widetilde{\scr{B}}^{\mathrm{inv}}$ 也是一个\trQuasiisomorphism，它诱导出%
\trDerivedCategory{}之间的一个\trEquivalence{}（\cref{thm:tw-quasi-equivalence-preserved}），
所欲证的结果随之而来。
\end{proof}
